\section{Compatible codewords and a change of self-dual code}
\label{sec:compatible-codes}

Put $N=4q$ with $q\ge3$ odd. A sign class means a pair $\{v,-v\}$;
use the representative whose first coordinate is positive. Set
$\pi(v)=\prod_jv_j$, which is well-defined on classes because $N$ is even.
We call a family modularly orthogonal if all its pairwise inner products
are divisible by $q$. The following proof requires no primality.

\begin{theorem}[Two-parity rigidity]\label{code:parity}
In a modularly orthogonal family of distinct sign classes, the
classes of each fixed parity are pairwise orthogonal over the reals.
If both parities occur, each parity has at most four classes.
Consequently any such family with at least nine classes is a real
partial Hadamard matrix, and every family has at most $N$ classes.
The mixed-parity bound of eight is sharp for every odd $q\ge3$.
\end{theorem}
\begin{proof}
For distinct nonantipodal signs $u,v$, their inner product is even,
divisible by odd $q$, and strictly between $-4q$ and $4q$. It is
therefore $0$ or $\pm2q$. Since
\[
 \pi(u)\pi(v)=(-1)^{(N-u\cdot v)/2},
\]
equal parities force zero and opposite parities force $\pm2q$.
Within one parity the normalized rows are orthonormal. Against
any normalized row of the opposite parity, every squared inner
product is $1/4$. Bessel's inequality limits the first class to
four rows, and then does the same for the second.
If only one parity occurs there are at most $N$ orthogonal rows;
otherwise there are at most eight, which is less than $N$.

For sharpness start with the four rows of $H_4$. Adjoin the four
antipodal classes of sign vectors of the form
$\frac12\sum_{i=1}^4\sigma_i h_i$ that have all entries $\pm1$.
The two sets are orthogonal bases of $\mathbb R^4$ whose cross
inner products are $\pm2$. Repeat each coordinate $q$ times.
This gives four classes of each parity, with cross products
$\pm2q$ and within-class products zero.
\end{proof}

This is an elementary rigidity result within the classical subject
of modular Hadamard matrices, not a priority claim; see Marrero and
Butson, \emph{Modular Hadamard matrices and related designs} (1973),
\url{https://doi.org/10.1016/0097-3165(73)90073-3}.
The modulus here grows with the target order. Existence theorems
for a fixed small modulus cannot simply supply this hypothesis.

\begin{corollary}[Five rows remove compatibility tests]\label{code:five}
Let $C$ be a submodule of $(\mathbb Z/q\mathbb Z)^N$ with
$C\subseteq C^\perp$, and let
\[
 S(C)=\{v\in\{\pm1\}^N:v\bmod q\in C,\ v_1=1\}.
\]
If $C$ contains five mutually orthogonal sign rows, every pair of
distinct members of $S(C)$ is orthogonal over the reals.
Thus $5\le |S(C)|\le N$.
\end{corollary}
\begin{proof}
The five prescribed orthogonal classes all have the same parity.
One opposite-parity class in $C$ would violate the four-row Bessel
bound. All sign classes in $C$ consequently have this same parity,
and Theorem~\ref{code:parity} applies.
\end{proof}

This changes the task from finding a clique among modular candidates
to counting sign words in one self-orthogonal code. It does not prove
that the count is large. In particular a dual kernel $P^\perp$ of a
short frame is usually not self-orthogonal; this corollary must not
be applied to that entire candidate set.
Five initial orthogonal rows do exist without assuming $H_N$:
write $N=12+8k$ and concatenate the first five rows of a fixed
$H_{12}$ with $k$ copies of the first five rows of $H_8$.
Their modular row span is self-orthogonal, but this construction
alone supplies no useful lower bound beyond five sign classes.

\subsection{A precise restricted-alphabet existence criterion}
The elementary three-column completion theorem proved in
\path{../round5/collision_entropy.tex} says that $N-3$ orthogonal
sign rows of length $N$ can be completed to a full Hadamard matrix.
For $N\ge12$, $N-3\ge9$. The preceding theorem therefore gives:
\[
 \begin{split}
 H_N\text{ exists}\quad\Longleftrightarrow\quad&
 \text{some self-orthogonal submodule }C\subseteq
 (\mathbb Z/q\mathbb Z)^N\\[-2pt]
 &\text{has }|S(C)|\ge N-3.
 \end{split}
\]
For the forward implication take the module generated by the rows
of $H_N$, whose inner products are zero modulo $q$, including the
diagonal $N$. For the reverse implication the sign words are a
partial Hadamard matrix by Theorem~\ref{code:parity}, and complete
its $N-3$ rows. This is an equivalent target, not a count estimate.
One may equivalently require the module to be self-dual. Indeed the
row module of a full $H_N$ is self-dual even for composite odd $q$:
if $Hz=qy$ for integer lifts, then $H^TH=4qI$ implies
$4z=H^Ty$. Since four is invertible modulo $q$, the residue of
$z$ belongs to the row module. Self-orthogonality gives the reverse
inclusion. This is the classical Hadamard-code construction, not
a new self-dual-code existence theorem; see Munemasa and Tamura,
\emph{The codes and the lattices of Hadamard matrices},
\url{https://doi.org/10.1016/j.ejc.2011.11.007}.

Only odd $q$ are needed for the full conjecture: write an arbitrary
multiple of four as $4q\,2^k$ with $q$ odd, and tensor an order-$4q$
matrix with a Sylvester matrix of order $2^k$. The case $q=1$
is already Sylvester. Neither this reduction nor the criterion
establishes existence at the remaining odd multipliers.

\subsection{A self-dual neighbor preserving existing sign rows}
Now assume $q=p$ is an odd prime and work over $\mathbb F_p$.
Let $C=C^\perp$ be self-dual, with at least five orthogonal sign
classes in $S(C)$. If a sign vector $v$ is orthogonal over the
integers to every member of $S(C)$ but does not belong to $C$, put
\[
 C'=(C\cap v^\perp)+\langle v\rangle.
\]
\begin{proposition}[Exact code pivot]\label{code:pivot}
The space $C'$ is self-dual, $\dim(C\cap C')=N/2-1$, and
$S(C)\subsetneq S(C')$. Conversely, if no real sign extension of
$S(C)$ exists, no self-orthogonal code preserving every member
of $S(C)$ can acquire a new sign class.
\end{proposition}
\begin{proof}
The functional $c\mapsto c\cdot v$ on $C$ is nonzero since
$C=C^\perp$ and $v\notin C$, so its kernel has dimension $N/2-1$.
Also $v\cdot v=N=0$ in $\mathbb F_p$. The displayed sum is
self-orthogonal of dimension $N/2$, hence self-dual.
It contains the old sign words because their real inner products
with $v$ vanish, and it contains the new class of $v$.
The assumption $v\notin C$ is automatic for a new real extension:
an extension is not one of the classes being extended, whereas
membership in $C$ would put its normalized representative in $S(C)$.
The converse follows from the same corollary: any acquired class
would be a real extension of all preserved classes.
\end{proof}

Thus repeated extension-guided pivots strictly increase the sign
count and stop after at most $N-|S(C)|$ moves, either at a full matrix
or at a genuinely terminal partial frame. This is not an algorithm
that escapes terminal frames. Moving to a different code can remove
an artificial restriction, but cannot make an incompatible new row
orthogonal to all existing rows.

\subsection{Exact known-order experiments}
The verifier exhausts all self-dual ternary codes containing the
first five rows of a specified Paley $H_{12}$. Their span has rank
five, and there are exactly two self-dual extensions, each with
twelve sign classes. These are results for that prescribed seed,
not a classification of all five-row frames.

At order twenty, a deterministic finite-field construction extends
the first five rows of a Paley $H_{20}$ to a self-dual code over
$\mathbb F_5$. Exhausting its $2^9$ normalized sign assignments in
systematic coordinates yields only eight classes. This disproves
the suggestion that any self-dual extension containing five rows
automatically gives a full matrix. The eight-row frame has rank
eight modulo five. Enumerating its modular kernel gives forty
classes: the eight existing rows and thirty-two exact extensions.
Using the first extension in a single pivot gives
\[
             |S(C)|=8\quad\longrightarrow\quad |S(C')|=20.
\]
All twenty rows are checked pairwise over the integers.
The finite-field bases and the extension witness are recorded in
\path{codes/checks.json}.

Finally self-duality alone need not give even one sign word.
Over $\mathbb F_3$, the length-four code generated by
$(1,0,1,1)$ and $(0,1,1,-1)$ is self-dual and has no word of full
support. Its threefold direct sum is a self-dual length-twelve code
with no sign words. A useful general lower bound must use more than
dimension and self-duality. No count near $N$ has been proved at
unknown orders.
